\section*{Source: RFE Response Brief.docx}
\textbf{Category}: source-7
\textbf{Page}: 1

\subsection*{Response to USCIS RFE – EB-1A I-140 \& I-485 Petition for Jessica Marie Thompson}

[Sponsor’s Full Name]
[Sponsor’s Address]
[City, State, ZIP Code]

Date: 2026-06-12

Subject: Response to Request for Evidence (RFE) – Jessica Marie Thompson

Dear USCIS Officer,

\subsection*{Introduction \& RFE Reference}

Parties \& Purpose:
“[Sponsor’s Full Name] (the “Petitioner”) submits this Response to the Request for Evidence (RFE) issued in connection with its EB-1A I-140 \& I-485 petition for Jessica Marie Thompson (the “Beneficiary”).” In this response, we address the concerns raised regarding the Beneficiary’s qualification under the extraordinary ability category. The petitioner has collated all available documentation as part of this petition. Every submission has been carefully reviewed and compiled to provide a complete record of the Beneficiary’s academic and professional background. We respectfully submit this detailed response for your consideration and further adjudication.

RFE Details:
Receipt Number: [Number]
RFE Issued: [Date]
Response Deadline: [Date]

\subsection*{Rebuttal to USCIS Concerns}

\subsection*{Concern 1: EB-1A Regulatory Criteria}

The USCIS has raised concerns regarding the Beneficiary’s ability to meet at least three of the ten regulatory criteria for EB-1A classification. In response, we outline multiple facets of the Beneficiary’s qualifications as demonstrated in our file. Though the record does not include all types of awards or memberships, the evidence provided – including a certified degree certificate from Veermata Jijabai Technological Institute showing academic excellence – supports her exceptional abilities in her field. Additional documentation such as the Form I-94 and visa pages further illustrate her history of lawful, exemplary performance. Each document contributes to a comprehensive demonstration of her sustained abilities and unique qualifications.

\subsection*{Concern 2: Sustained National or International Acclaim}

USCIS has questioned whether the evidence demonstrates sustained national or international acclaim. In rebuttal, we emphasize that the Beneficiary’s academic record, as evidenced by her prestigious degree certificate from Veermata Jijabai Technological Institute, clearly reflects a history of outstanding achievement. Moreover, the supportive records including the Form I-94, visa pages, and passport validate her continuous travel and engagement in international academic and professional circles. These documents, while not exhaustive of every accolade, collectively underscore her recognized status in her area of expertise. Through these submissions, we provide a clear timeline of achievements, reinforcing her long-standing and internationally acknowledged contributions.

\subsection*{Concern 3: Intent to Continue Work in the Area of Extraordinary Ability}

The USCIS further expressed concerns regarding the Beneficiary’s intent to continue working within her area of extraordinary ability after petition approval. While the file does not include explicit employment offer letters or contracts from an employer, the Beneficiary’s academic achievements and documented legal admissions, as seen in her I-94 record and visa documentation, support her deep commitment to further contributions in her field. Her academic credentials and sustained travel records reflect a professional trajectory aimed at continuous improvement and engagement. These exhibits confirm that she remains fully prepared and intent on making significant future contributions, further establishing her commitment to excellence in her specialty.

\subsection*{Additional Legal Authority}

In support of our petition, we respectfully reference INA § 203(b)(1)(A) and 8 C.F.R. § 204.5(h), which outline the standard for classifying aliens possessing extraordinary ability. Throughout our submission, we have systematically addressed the relevant regulatory criteria with supporting documentation including, but not limited to, the Beneficiary’s degree certificate, Form I-94, visa pages, and passport. This legal framework substantiates the petitioner’s assertion that the Beneficiary has demonstrated the required sustained acclaim and will continue to contribute at a high level within her field. We trust that these clarifications and references to established legal authority affirm the strength of our petition.

\subsection*{Conclusion \& Request}

Based on the evidence submitted and the legal authorities cited, the Petitioner respectfully reaffirms that the Beneficiary, Jessica Marie Thompson, meets the requisite criteria for classification as an alien of extraordinary ability under INA § 203(b)(1)(A). The enclosed documents – including the degree certificate, Form I-94, visa pages, and passport – collectively underscore her proven record and commitment to further contributions in her field. Petitioner respectfully requests prompt adjudication of the EB-1A I-140 \& I-485 petition and kindly asks for notification via email at [Email Address]. For additional inquiries or documentation requests, please contact [Sponsor’s Representative Name], [Title], at [Phone Number] or [Email Address].

Very truly yours,
\_\_\_\_\_\_\_\_\_\_\_,
[Sponsor’s Representative Name], [Title]
[Sponsor’s Organization Name]